\section{Related Work}

\noindent\textbf{Knowledge Graph Retrieval-Augmented Generation (KG-RAG).}
KG-RAG augments LLMs with structured KGs to improve quality and compositional reasoning. Early work grounds LLMs' generation by translating natural-language questions into executable graph queries (e.g., SPARQL/Cypher), retrieving relevant subgraphs or answers, and feeding them back to the model~\citep{ouyang2022traininglanguagemodelsfollow, lee2025sparkle}. More recent approaches adopt a modular LLM pipeline over KGs, interleaving natural language reasoning with multi-stage planning, path search, and execution, where each stage uses a prompted or fine-tuned LLM~\citep{Luo2024RoG, Sun2024ToG, Chen2024Pog, Wang2025ReKnoS}. Despite these advances, most systems rely on hand-engineered modules or prompt heuristics tied to a specific KG schema, which induce cost inefficiency and limit generalization. These challenges motivate our \framework{} framework: a single-agent KG-RAG approach that improves efficiency and transfers well to new KGs.

\noindent\textbf{Multi-turn RL for LLMs.}
Reinforcement learning (RL) has become central to equipping LLMs with step-by-step (chain-of-thought) reasoning behavior~\citep{OpenAI2024o1, DeepSeek2025R1}. RL-enhanced models yield substantial gains in math and coding~\citep{le2022coderl,chervonyi2025goldmedal}, and broader complex reasoning tasks.
More recently, RL has been applied to agentic LLMs that invoke external tools (e.g., bash terminals and APIs) or interact with knowledge bases, improving tool use~\citep{Qin2023ToolLLM} and retrieval-augmented generation (RAG)~\citep{jin2025searchr1,wang2025otc} to facilitate tool use or RAG. Building on these advances, our work, \framework{} adopts end-to-end RL as the core optimization algorithm for an agentic KG-RAG framework.


